\ifdefined\iscombined
	\lecturetitle{Op-amps and Active Circuits}
\else
\documentclass{article}
\usepackage{xparse}
\usepackage{changepage}
\usepackage{listings}
\usepackage{amsthm}
\usepackage{comment}
\usepackage{amsmath}
\usepackage{braket}
\usepackage{amsfonts}
\usepackage{sectsty}
\usepackage{hyperref}
\usepackage{fancyhdr}
\usepackage[top=1in,bottom=1in,left=1in,right=1in]{geometry}
\usepackage{float}
\usepackage{graphicx}
\usepackage{mathtools}
\usepackage{tikz}
\usetikzlibrary{arrows}
\usepackage{circuitikz}

\date{
	\vspace{-.75em}
	{\large \today}
}

\title{
	\vspace*{-1.5em}
	{\Huge Lecture 8} \\
	\vspace{-1em}
	\rule{\textwidth/4}{.01em} \\
	\vspace{-.25em}
	{\Large Op-amps and Active Circuits}
	\vspace{-.75em}
}

\author{
	{\large Matthew Farah}
}

\hypersetup{
	colorlinks=true,
	linkcolor=blue,
	filecolor=magenta,
	urlcolor=blue,
	citecolor=blue,
	anchorcolor=blue
}

\fancyhead{}
\fancyhead[L]{\today}
\fancyhead[R]{Lecture 8 - Op-amps and Active Circuits}

\setlength{\parindent}{0cm}

\newcounter{example}
\renewcommand{\theexample}{\arabic{example}}

\newenvironment{example}[1][]{
	\refstepcounter{example}
	\begin{adjustwidth}{1.5em}{}
	\noindent\textbf{\ifx&#1&Example \theexample:\else Example \theexample \ (#1):\fi}
	\ignorespaces
}{
	\vspace{.5em}
	\end{adjustwidth}
}

\newenvironment{solution}{
	\vspace{.5em}
	\par
	\begin{adjustwidth}{1.5em}{}
	\textbf{{Solution:}}
	\par
	\vspace{.5em}
}{
	\vspace{1em}
	\end{adjustwidth}
}

\newcounter{subexample}
\renewcommand{\thesubexample}{\alph{subexample}}

\newenvironment{subexample}{
	\setcounter{step}{0}
	\refstepcounter{subexample}
	\vspace{1em}\par\noindent
	\begin{adjustwidth}{1.5em}{}
	\noindent\textbf{(\thesubexample)}
	\ignorespaces
}{
	\par
	\vspace{1em}
	\end{adjustwidth}
}

\renewenvironment{proof}{
	\vspace{.5em}\par\noindent
	\begin{adjustwidth}{1.5em}{}
	\emph{Proof:}\par\vspace{1em}
	\setlength{\parindent}{0cm}
}{
	\end{adjustwidth}
}

\newtheorem{theorem}{Theorem}

\renewenvironment{theorem}[1][]{
	\refstepcounter{theorem}
	\vspace{1em}\par\noindent
	\begin{adjustwidth}{1.5em}{}
	\noindent\textbf{\ifx&#1&Theorem \thetheorem:\else Theorem \thetheorem \ (#1):\fi}
	\ignorespaces
}{
	\vspace{.5em}
	\end{adjustwidth}
}

\newtheorem{lemma}{Lemma}

\renewenvironment{lemma}[1][]{
	\refstepcounter{lemma}
	\vspace{1em}\par\noindent
	\begin{adjustwidth}{1.5em}{}
	\noindent\textbf{\ifx&#1&Lemma \thelemma:\else Lemma \thelemma \ (#1):\fi}
	\ignorespaces
}{
	\vspace{.5em}
	\end{adjustwidth}
}

\makeatother
\makeatletter

\newcounter{step}
\renewcommand{\thestep}{\Roman{step}}

\newenvironment{step}{
	\refstepcounter{step}
	\vspace{1em}\par\noindent
	\ignorespaces\noindent\textbf{(\thestep)}
	\ignorespaces
}{
	\par
}

\sectionfont{\fontsize{12}{12}\bfseries}
\subsectionfont{\fontsize{10}{12}\normalfont}
\subsubsectionfont{\fontsize{10}{12}\normalfont}

\begin{document}

\maketitle
\pagestyle{fancy}

\fi

\begin{itemize}
	\item Quiz next week will be on circuit analysis.
	\item Op-amps probably won't be on the midterm.
	\item Operational amplifiers (op-amps) are a kind of circuit element used to produce `active' circuits.
	\begin{itemize}
		\item Active circuits are used to implement transfer functions as a circuit.
	\end{itemize}
	\item Passive circuits only use passive circuit elements.
	\begin{itemize}
		\item Passive circuit elements only consume energy.
		\begin{itemize}
			\item Ex. Resistors, capacitors, etc.
		\end{itemize}
		\item All circuits seen thus far in the course used only passive elements (besides the source voltage).
	\end{itemize}
	\item Active circuits use one or more active circuit elements.
	\begin{itemize}
		\item Active circuit elements amplify, produce, or switch energy.
		\begin{itemize}
			\item Ex. Source voltages, op-amps, etc.
		\end{itemize}
	\end{itemize}
	\item Op-amps effectively work like low-pass filters with a very high cut-off.
	\item Op-amps look like this:
\end{itemize}

\vspace{1em}

\begin{figure}[H]
	\centering
	\begin{circuitikz}
		\draw (0,0) node[op amp] (opamp) {$A$};
		\draw (opamp.-) -- ++(-1,0) node[left]{$+V_1(t)$};
		\draw (opamp.+) -- ++(-1,0) node[left]{$+V_2(t)$};
		\draw (opamp.out) -- ++(1,0) node[right]{$V_o(t)$};
		\node[above=2pt] at ([xshift=-5pt]opamp.-) {$\scriptstyle i_-$};
		\node[below=2pt] at ([xshift=-5pt]opamp.+) {$\scriptstyle i_+$};
	\end{circuitikz}
\end{figure}

\vspace{1em}

Where:

\begin{itemize}
	\item $V_1(t)$ and $V_2(t)$ are its input voltages.
	\item $i_-(t)$ and $i_+(t)$ are its input currents.
	\item $V_o(t)$ is its output voltage.
	\item $A$ is its gain.
\end{itemize}

\vspace{.5em}

\begin{itemize}
	\item Op-amps have four important characteristics in an ideal setting:
	\begin{enumerate}
		\item $V_o(t) = A(V_2(t) - V_1(t))$.
		\item High input impedance.
		\begin{itemize}
			\item I.E. $Z_i = \infty$.
		\end{itemize}
		\item Low output impedance.
		\begin{itemize}
			\item I.E. $Z_o = 0$.
		\end{itemize}
		\item High gain constant.
		\begin{itemize}
			\item I.E. $A = \infty$.
		\end{itemize}
	\end{enumerate}
	\item Ideal op-amps also:
	\begin{itemize}
		\item Have equal input voltages.
		\begin{itemize}
			\item $V_1(t) = V_2(t)$.
		\end{itemize}
		\item Have zero input currents.
		\begin{itemize}
			\item $i_-(t) = i_+(t) = 0$.
		\end{itemize}
	\end{itemize}
	\item The relationship between the input and output voltages of an op-amp is linear around a specific input voltage.
	\begin{itemize}
		\item The saturation voltages are the voltages where this is no longer true.
	\end{itemize}
	\item Chaining op-amps together makes them more stable.
	\begin{itemize}
		\item \emph{Aside:} Not too sure what exactly that means.
	\end{itemize}
\end{itemize}

\ifdefined\iscombined
	\newpage
\else
	\end{document}
\fi
